\section{The Parry-rank fold for every quadratic Pisot base}
\label{sec:quadratic-pisot-extension}

This section closes the interface left by Remark~\ref{rem:metallic-beta-scope}.
The point requiring care is that a quadratic Pisot number need not be a unit
and its conjugate need not be negative.  Consequently the metallic
Ostrowski language is not the appropriate common object.  We instead use the
greedy $\beta$-language itself and its canonical colexicographic rank.  The
published results on greedy admissibility and finite-state normalization are
used only for context; the finite rank and sliding-congruence arguments needed
below are proved here.

\subsection{Quadratic classification and canonical finite windows}

\begin{lemma}[The two quadratic Pisot parameter chambers]
\label{lem:quadratic-pisot-chambers}
Let $\beta>1$ be a quadratic algebraic integer and let $\beta'$ be its
conjugate.  Then $|\beta'|<1$ if and only if the minimal polynomial of
$\beta$ belongs to exactly one of the following two families:
\begin{align*}
P^-_{a,b}(x)&=x^2-ax-b, & a\ge b\ge1,\\
P^+_{a,b}(x)&=x^2-ax+b, & a\ge3,\quad 1\le b\le a-2.
\end{align*}
In the first chamber $\beta'<0$, $\lfloor\beta\rfloor=a$, and
\[
d_\beta(1)=ab0^\infty,\qquad
d_\beta^*(1)=\bigl(a(b-1)\bigr)^\infty.
\]
In the second chamber $\beta'>0$, $\lfloor\beta\rfloor=a-1$, and
\[
d_\beta(1)=(a-1)(a-b-1)^\infty=d_\beta^*(1).
\]
Here $d_\beta$ is the greedy expansion and $d_\beta^*$ is Parry's
quasi-greedy boundary word.
\end{lemma}

\begin{proof}
Write the monic minimal polynomial as $x^2-ax+c$.  Since
$c=\beta\beta'\ne0$, the sign of $c$ is the sign of $\beta'$.  If
$c=-b<0$, then $b/\beta=|\beta'|<1$.  The positive root exceeds $a$;
the last inequality is equivalent, for integral $a,b$, to $1\le b\le a$.
Conversely these inequalities place the negative root in $(-1,0)$.
If $c=b>0$, then $0<b/\beta=\beta'<1$ and $a-1<\beta<a$; evaluating
$P^+_{a,b}$ at $a-1$ and $a$ gives exactly $a\ge3$ and
$1\le b\le a-2$.  The converses follow from the same sign test.

For $P^-_{a,b}$ one has $\beta^2=a\beta+b$.  The greedy algorithm applied
to $1$ therefore gives the two digits $a,b$ and then terminates; Parry's
standard replacement of a finite expansion gives
$d_\beta^*(1)=(a(b-1))^\infty$.  For $P^+_{a,b}$, put
$r=\beta-a+1=1-b/\beta$.  Then
$\lfloor\beta\rfloor=a-1$ and
\[
\beta r=\beta-b=(a-b-1)+r,
\]
so every digit after the first is $a-b-1$.  The expansion is infinite and
hence already quasi-greedy.  This is precisely Parry's calculation
\cite{Parry1960}; it is included to fix the parameters, not as a new proof of
his general admissibility theorem.
\end{proof}

For either chamber put $d=\lfloor\beta\rfloor$ and
$\mathcal A_\beta=\{0,\ldots,d\}$.  Let $\mathcal L_m(X_\beta)$ be the
length-$m$ language of the two-sided greedy $\beta$-shift, characterized by
Parry's suffix inequalities against $d_\beta^*(1)$.

\begin{definition}[Canonical Parry-rank window fold]
\label{def:quadratic-parry-fold}
For $m\ge0$ define $Q_m(\beta)=|\mathcal L_m(X_\beta)|$.  In the two
chambers of Lemma~\ref{lem:quadratic-pisot-chambers}, respectively, these
integers are
\begin{align}
Q_0&=1,\quad Q_1=a+1,\quad Q_{j+2}=aQ_{j+1}+bQ_j,
&& (P^-_{a,b}),\label{eq:qminus}\\
Q_0&=1,\quad Q_1=a,\quad Q_{j+2}=aQ_{j+1}-bQ_j,
&& (P^+_{a,b}).\label{eq:qplus}
\end{align}
Use the low-to-high orientation
\[
X_m^{(\beta)}:=\{x\in\mathcal A_\beta^m:
(x_m,\ldots,x_1)\in\mathcal L_m(X_\beta)\}
\]
and set
\[
\operatorname{Rank}_{\beta,m}(x)=\sum_{j=0}^{m-1}x_{j+1}Q_j(\beta).
\]
The canonical finite-window fold is
\[
\Fold_{\beta,m}(w)=
\bigl(\operatorname{Rank}_{\beta,m}|_{X_m^{(\beta)}}\bigr)^{-1}
\left(\operatorname{Rank}_{\beta,m}(w)\bmod Q_m(\beta)\right).
\]
Finally,
\[
(\Phi_{\beta,m}(z))_t
=\Fold_{\beta,m}(z_t,\ldots,z_{t+m-1}),\qquad
Y_{\beta,m}=\Phi_{\beta,m}(\mathcal A_\beta^{\mathbb Z}).
\]
The nontrivial threshold is
\[
m_*(\beta)=\min\{m\ge2:\Phi_{\beta,m}\text{ is a conjugacy onto the
SFT }Y_{\beta,m}\}.
\]
\end{definition}

\begin{proposition}[Parry rank is consecutive and canonical]
\label{prop:parry-rank-bijection}
For every quadratic Pisot $\beta$ and $m\ge1$,
\[
\operatorname{Rank}_{\beta,m}|_{X_m^{(\beta)}}:
X_m^{(\beta)}\longrightarrow\{0,\ldots,Q_m(\beta)-1\}
\]
is an order-preserving bijection when $X_m^{(\beta)}$ is read in the
colexicographic order.  Thus Definition~\ref{def:quadratic-parry-fold} is
determined by the greedy language and its Parry word; it is not a choice of
an auxiliary enumeration.  For $P^-_{a,1}$ it is exactly
Definition~\ref{def:metallic-window-data}.
\end{proposition}

\begin{proof}
In the negative chamber, Parry admissibility reduces in the low-to-high
orientation to
\[
x_k=a\Longrightarrow x_{k-1}<b.
\]
Words whose last digit is $0,\ldots,a-1$ give $a$ consecutive intervals of
length $Q_{m-1}$.  If the last digit is $a$, the preceding digit is one of
$0,\ldots,b-1$ and gives $b$ consecutive intervals of length $Q_{m-2}$.
Their union is
$[0,aQ_{m-1}+bQ_{m-2}-1]$.  This proves both \eqref{eq:qminus} and the rank
claim by induction.

In the positive chamber put $c=a-b-1$.  Parry admissibility says that after
an occurrence of $a-1$ the remaining lower word is at most $c^\infty$ in
lexicographic order.  Simultaneous induction gives
\[
|\mathcal L_m(X_\beta)|=aQ_{m-1}-bQ_{m-2}
\]
and says that the legal lower words bounded by $c^\infty$ have ranks
\[
0,\ldots,Q_{m-1}-bQ_{m-2}-1.
\]
Indeed the last digit classes $0,\ldots,a-2$ contribute the consecutive
interval $[0,(a-1)Q_{m-1}-1]$, while the last digit $a-1$ contributes the
next interval of the displayed bounded length.  Splitting the bounded class
according as its next digit is smaller than or equal to $c$ gives the same
statement one index lower, which initializes at lengths zero and one.
The upper endpoint is
$(a-1)Q_{m-1}+Q_{m-1}-bQ_{m-2}-1=Q_m-1$.
This proves \eqref{eq:qplus}, consecutiveness, and order preservation.
When $b=1$ in the negative chamber, the rule and recurrence are precisely
those of Definition~\ref{def:metallic-window-data}.
\end{proof}

\subsection{The bounded sliding-congruence lemma}

The following elementary lemma is the arithmetic core.  It is stated
separately so that the exact hypothesis behind the family-wide claim is
visible.

\begin{lemma}[Quadratic bounded sliding congruences]
\label{lem:quadratic-sliding-congruences}
Take either parameter chamber in
Lemma~\ref{lem:quadratic-pisot-chambers}, let $d=\lfloor\beta\rfloor$, and
let $(Q_j)$ be \eqref{eq:qminus} or \eqref{eq:qplus}.  If $m\ge3$ and
$e_0,\ldots,e_{2m-2}\in[-d,d]\cap\mathbb Z$ satisfy
\begin{equation}
\sum_{j=0}^{m-1}Q_j e_{t+j}\equiv0\pmod{Q_m}
\qquad(0\le t<m),\label{eq:sliding-kernel}
\end{equation}
then $e_0=\cdots=e_{2m-2}=0$.
\end{lemma}

\begin{proof}
Write $s=1$ in the negative chamber and $s=-1$ in the positive chamber, so
$Q_{j+2}=aQ_{j+1}+sbQ_j$, and put
$\delta=Q_1-a$; hence $(d,\delta)=(a,1)$ for $s=1$ and
$(d,\delta)=(a-1,0)$ for $s=-1$.  Subtracting $a$ times the congruence at
$t+1$ and $sb$ times that at $t+2$ from the congruence at $t$ cancels every
interior coefficient and gives
\begin{equation}
e_t+\delta e_{t+1}-sbQ_{m-1}e_{t+m+1}\equiv0\pmod{Q_m}.
\label{eq:quadratic-elimination}
\end{equation}

We record the elementary separation calculation used twice below.  Direct
division using the recurrence and the parameter inequalities gives
\begin{equation}
\operatorname{dist}(beQ_{r-1},Q_r\mathbb Z)\ge Q_{r-2}
\quad(r\ge4,\ 1\le |e|\le d).
\label{eq:quadratic-separation}
\end{equation}
For completeness, the division has only the following two rows; replacing
$e$ by $-e$ changes no distance.
\[
\begin{array}{c|c|c|c}
s&Q_1&\text{parameter range}&Q_2\text{ and required lower bound}\\ \hline
+1&a+1&1\le b\le a&a^2+a+b>2a\\
-1&a&1\le b\le a-2&a^2-b>a-1
\end{array}
\]
To verify \eqref{eq:quadratic-separation}, choose the nearest integer $k$
and insert $Q_r=aQ_{r-1}+sbQ_{r-2}$ in
$beQ_{r-1}-kQ_r$.  If its modulus were below $Q_{r-2}$, division by
$Q_{r-2}$ and the inequalities
\[
a<\frac{Q_{r-1}}{Q_{r-2}}<a+b\quad(s=1),
\qquad
a-1<\frac{Q_{r-1}}{Q_{r-2}}<a\quad(s=-1)
\]
force respectively
$be=ka+ksbQ_{r-2}/Q_{r-1}$ with a nonzero integer strictly between
$-1$ and $1$.  This is impossible.  Equality can occur only at an endpoint
and is harmless.  This proves the separation claim by induction from the
displayed $Q_2$ row.

Suppose first that $m\ge4$.  The small part of
\eqref{eq:quadratic-elimination} has modulus at most
$d(1+\delta)$, whereas \eqref{eq:quadratic-separation} and the table give a
strictly larger distance whenever $e_{t+m+1}\ne0$.  Hence
$e_{t+m+1}=0$ for $0\le t\le m-3$.  Thus
$e_{m+1}=\cdots=e_{2m-2}=0$.  The final original congruence reduces to
\[
e_{m-1}+Q_1e_m\equiv0\pmod{Q_m}.
\]
Its left side has modulus at most $d(1+Q_1)<Q_m$ and $Q_1>d$, so both
digits vanish.  Back substitution in the preceding congruences then gives
$e_{m-2},\ldots,e_0=0$ successively.

It remains to check $m=3$, where endpoint equality in the separation bound
can occur.  Here each integer sum in \eqref{eq:sliding-kernel} has modulus
less than $2Q_3$, so it is $k_tQ_3$ with $k_t\in\{-1,0,1\}$.  Substituting
the three values of $t$ gives
\begin{equation}
e_t+Q_1e_{t+1}+Q_2e_{t+2}=k_tQ_3,
\qquad t=0,1,2.
\label{eq:quadratic-m3-equalities}
\end{equation}
We now eliminate the two conjugate signs separately.

Suppose first that $s=1$.  Thus $d=a$ and
\[
(Q_1,Q_2,Q_3)
=(a+1,a^2+a+b,a^3+a^2+2ab+b).
\]
Consider one equation $x+Q_1y+Q_2z=kQ_3$ occurring in
\eqref{eq:quadratic-m3-equalities}.  By symmetry it suffices to treat
$k=0,1$.  If $k=0$, then $|z|\ge2$ is impossible because
$2Q_2>a(1+Q_1)$.  The choice $z=0$ gives $x=y=0$, while $z=1$ gives
\[
x+b+(a+1)(y+a)=0.
\]
Since $y+a\ge0$ and $-(x+b)\le a-b<a+1$, this forces
$(x,y,z)=(-b,-a,1)$; the case $z=-1$ is its negative.  If $k=1$ and
$z\le a-1$, then
\[
x+Q_1y+Q_2z
\le a+aQ_1+(a-1)Q_2
=Q_3-\bigl(b(a+2)-a\bigr)<Q_3.
\]
Hence $z=a$, after which $x+(a+1)y=b(a+1)$ forces $(x,y)=(0,b)$.
Consequently every bounded solution of one equation is
\[
(x,y,z;k)\in
\{(0,0,0;0),\ \pm(b,a,-1;0),\ \pm(0,b,a;1)\}.
\]
The terminal pairs of these triples are
$(0,0)$, $\pm(a,-1)$, and $\pm(b,a)$.  If $e_4\ne0$ and the $t=2$
triple is $\pm(0,b,a)$, its initial pair $\pm(0,b)$ cannot be a terminal
pair of the $t=1$ triple.  If instead the $t=2$ triple is
$\pm(b,a,-1)$, the overlap forces the $t=1$ triple to be
$\pm(0,b,a)$, whose initial pair $\pm(0,b)$ cannot be a terminal pair of
the $t=0$ triple.  Both alternatives are impossible, and therefore
$e_4=0$.

Suppose next that $s=-1$.  Put $c=a-b$ and $d=a-1$; then
$2\le c\le d$ and
\[
(Q_1,Q_2,Q_3)=(a,a^2-b,a^3-2ab).
\]
For one equation $x+ay+Q_2z=kQ_3$, again it suffices to take $k=0,1$.
When $k=0$, the inequality $2Q_2>d(1+a)=a^2-1$ gives $|z|\le1$.
The choice $z=0$ gives $x=y=0$, whereas $z=1$ gives
\[
x-b+a(y+a)=0.
\]
Here $1\le y+a\le2a-1$ and $b-x\le b+a-1\le2a-3$, so necessarily
$y+a=1$ and $x=b-a=-c$; the case $z=-1$ is the negative solution.
For $k=1$, if $z\le a-2$, then
\[
x+ay+Q_2z
\le d(1+a)+(a-2)Q_2
=Q_3-\bigl(a(a-b)-2b+1\bigr)<Q_3,
\]
because $a-b\ge2$.  Thus $z=d$, and
\[
x+ay=a(a-b)-b=ac-b.
\]
Reduction modulo $a$, together with $|x|\le a-1$, leaves precisely
$(x,y)=(-b,c)$ and $(x,y)=(c,c-1)$.  Hence the complete list is
\[
(x,y,z;k)\in
\{(0,0,0;0),\ \pm(c,d,-1;0),\
  \ \pm(-b,c,d;1),\ \pm(c,c-1,d;1)\}.
\]
Its terminal pairs are $(0,0)$, $\pm(d,-1)$, $\pm(c,d)$, and
$\pm(c-1,d)$.  Since $2\le c\le d$, neither $\pm(-b,c)$ nor
$\pm(c,c-1)$ occurs in this list of terminal pairs.  Thus a $t=2$ triple
of either of the last two types cannot overlap a $t=1$ triple.  The only
remaining nonzero possibility is $\pm(c,d,-1)$; it forces the $t=1$
triple to be $\pm(-b,c,d)$, whose initial pair $\pm(-b,c)$ cannot overlap
a $t=0$ triple.  Therefore $e_4=0$ also in this chamber.

In either sign, every nonzero triple in the corresponding list has nonzero
last coordinate.  It follows successively from the equations at
$t=2,1,0$ that $e_2=e_3=e_4=0$, then $e_1=0$, and finally $e_0=0$.
Thus the kernel is trivial also at $m=3$.  This endpoint calculation is
independently certified by the routine \texttt{block\_language\_profile}
in \url{artifacts/verify_quadratic_pisot_threshold.py}, called with
$m=n=3$ in its threshold-and-block-language battery.
\end{proof}

\subsection{Sharp threshold and extremal bases}

\begin{theorem}[Full quadratic-Pisot fold threshold]
\label{thm:full-quadratic-pisot-threshold}
Let $\beta$ be an arbitrary quadratic Pisot number and use exactly the
greedy language, Parry data, rank, cyclic modulus and fold of
Definition~\ref{def:quadratic-parry-fold}.  Then
\[
m_*(\beta)=
\begin{cases}
3,&\beta\text{ is the root of }x^2-ax-a,\quad a\ge1,\\
3,&\beta\text{ is the root of }x^2-ax+1,\quad a\ge3,\\
2,&\text{otherwise}.
\end{cases}
\]
More strongly, for every $m\ge m_*(\beta)$ and every $n\ge m$, the
finite-block sliding fold is injective and
\begin{equation}
|\mathcal L_n(Y_{\beta,m})|
=(\lfloor\beta\rfloor+1)^{n+m-1}.
\label{eq:quadratic-pisot-block-count}
\end{equation}
At these apertures $\Phi_{\beta,m}$ is a topological conjugacy from the full
$(\lfloor\beta\rfloor+1)$-shift onto the SFT $Y_{\beta,m}$, with inverse of
memory zero and anticipation $m-1$.

Consequently the maximum of the threshold function on the full quadratic
Pisot family is three.  Its extremizers are exactly the negative-conjugate
diagonal $P^-_{a,a}$ and the positive-conjugate norm-one family
$P^+_{a,1}$.  The golden mean is the least extremizer, but it is not the
unique extremizer after passage from metallic units to all quadratic Pisot
bases.
\end{theorem}

\begin{proof}
Equality of two folded output blocks is, by
Proposition~\ref{prop:parry-rank-bijection}, equivalent to the sliding
congruences \eqref{eq:sliding-kernel} for their digit differences.  For
$m\ge3$, Lemma~\ref{lem:quadratic-sliding-congruences} makes every
length-$m$ finite-block map injective.  Consecutive decoded blocks cover a
longer raw block, proving injectivity for every $n\ge m$ and the count
\eqref{eq:quadratic-pisot-block-count}.  The standard gluing argument used in
Theorem~\ref{thm:metallic-threshold-classification} then gives a local inverse
and an SFT description by forbidden output words of length at most $m+1$.

It remains to classify $m=2$.  Put
\[
\kappa=Q_2-dQ_1=
\begin{cases}b,&P^-_{a,b},\\ a-b,&P^+_{a,b}.
\end{cases}
\]
Proposition~\ref{prop:parry-rank-bijection} gives the same explicit rule in
both chambers:
\[
\Fold_{\beta,2}(u,v)=
\begin{cases}
(u,v),&v<d\text{ or }u<\kappa,\\
(u-\kappa,0),&v=d\text{ and }u\ge\kappa.
\end{cases}
\]
If the second coordinate of the preceding output is positive, it is the raw
central digit.  If it is zero, that digit is either $0$ or $d$; the first
coordinate of the next output is respectively $0$ or one of
$d-\kappa,d$.  Hence the raw digit is locally recoverable exactly when
$d-\kappa>0$.  This condition is $a>b$ in the negative chamber and $b>1$
in the positive chamber.  If $d=\kappa$, the two constant raw points
$0^{\mathbb Z}$ and $d^{\mathbb Z}$ both map to the constant zero output,
so conjugacy fails.  Thus failure occurs exactly for $P^-_{a,a}$ and
$P^+_{a,1}$, and Lemma~\ref{lem:quadratic-sliding-congruences} shows that the
next aperture succeeds.  The chamber classification is exhaustive by
Lemma~\ref{lem:quadratic-pisot-chambers}, completing the proof.
\end{proof}

\begin{remark}[What is and is not claimed]
\label{rem:quadratic-open-interface}
The theorem concerns the canonical cyclic Parry-rank fold of
Definition~\ref{def:quadratic-parry-fold}.  Published finite-transducer
normalization results for Pisot bases
\cite{Frougny1992,FrougnySteiner2008,BertheFrougnyRigoSakarovitch2019}
establish recognizability for other redundant digit sets and boundary
conventions, but they do not identify those folds with this cyclic rank map.
No threshold claim is made for signed digits, non-greedy expansions,
non-cyclic truncation, degree greater than two, or general Pisot substitution
codings.  These choices are genuine residual interfaces rather than
consequences of Theorem~\ref{thm:full-quadratic-pisot-threshold}.
\end{remark}
